% !TeX root = ../tesis.tex
%%%%%%%%%%%%%%%%%%%%%%%%%%%%%%%%%%%%%%%%%%%%%%%%%%%%%%%%%%%%%%%%%%%
% SECCIÓN 6.2 — ESCENARIOS Y PROCESAMIENTO
% Issue #25. Pocas figuras: los datos van en cuadros; los mapas de
% los escenarios se usan en §6.3 (indicadores).
%
% PENDIENTE (decisión del autor, 2026-10-18): falta una sección
% específica que defina los dos escenarios (Caso 2 = Anillo MB,
% Caso 3 = Anillo Metro) y los explique con sus mapas. Se desarrolla
% al terminar este archivo; no antes.
%
% PENDIENTE para la revisión final (#32): Cap. 3 (3-3-5) y Cap. 5
% (5-4-2) dicen que T_b depende del intervalo de paso de cada
% sistema; VFTModel usa un valor fijo (FALLBACK_FRECUENCIA/2). Por
% eso aquí no se dan los intervalos de los anillos (3 y 5 min).
%%%%%%%%%%%%%%%%%%%%%%%%%%%%%%%%%%%%%%%%%%%%%%%%%%%%%%%%%%%%%%%%%%%
\section{Escenarios y Procesamiento}
\label{sec:anillo-escenarios}

Los tres casos se calculan con el mismo procedimiento del
Capítulo~\ref{ch:analisis-red-actual}. Lo único que cambia de uno a otro es el
\textsc{api-i} y la forma en que opera: el Caso~1 no lo tiene, el Caso~2 lo
incorpora con la operación del Metrobús y el Caso~3 con la del Metro. El código,
la corrida y las reglas con las que se arma el grafo son las mismas en los tres;
en particular, los transbordos peatonales se generan entre estaciones a 85~m o
menos (§\ref{sec:grafo-topologico}) y la penalización por transbordo se aplica
como en §\ref{subsec:fase2-tb}. Así, cualquier diferencia en los indicadores se
puede atribuir al anillo y no a un cambio de método. Esta sección reúne los
parámetros de cada caso, resume cómo se procesaron con la cadena de
herramientas y describe el tamaño del grafo que resulta.

\subsection{Parámetros de los Casos}
\label{subsec:anillo-parametros}

Cada caso toma los parámetros de un sistema que ya opera en la red, para que
el anillo se calcule con las mismas reglas que el resto del grafo. El Caso~3
replica al Metro: 36~km/h sobre vía exclusiva, sin fricción, igual que las
líneas del Metro en el Cuadro~\ref{tab:validacion-impedancia}. El Caso~2 parte
de los 16.3~km/h del Metrobús, pero en \texttt{Apímetro} se registró con
derecho de vía confinado, la categoría del Mexibús, y no exclusivo como las
líneas de Metrobús de la red actual; con $CF = 1.152$ su velocidad efectiva
baja a unos 14.2~km/h, la misma que alcanza el Mexibús. Fuera de la operación,
los dos casos son idénticos: comparten el trazo, las 94 estaciones y los ocho
ramales descritos en §\ref{subsec:corredor-anillar-cuadrantes}. El
Cuadro~\ref{tab:anillo-parametros} reúne estos valores.

\begin{table}[H]
    \centering
    \caption[Parámetros de los casos con anillo]{%
    Parámetros de operación del \textsc{api-i} en los Casos~2 y~3.
    Corrida: 2026-10-04.}
    \label{tab:anillo-parametros}
    \small
    \begin{tabular}{|l|c|c|}
        \hline
        \textbf{Parámetro} & \textbf{Caso~2 (Anillo MB)} &
        \textbf{Caso~3 (Anillo Metro)} \\
        \hline
        Sistema que replica       & Metrobús   & Metro \\
        Derecho de vía            & Confinado  & Exclusivo \\
        $\alpha$                  & 0.2        & 0.0 \\
        $CF$                      & 1.152      & 1.000 \\
        $V_{\text{base}}$ (km/h)  & 16.3       & 36.0 \\
        $V$ efectiva (km/h)       & 14.2       & 36.0 \\
        Estaciones                & 94         & 94 \\
        Líneas / ramales          & 4 / 8      & 4 / 8 \\
        \hline
    \end{tabular}
    \fuentefigura{Elaboración propia con \texttt{Apímetro} y
    \texttt{VFTModel}~\citep{cabrera2025apimetro, cabrera2026VFTModel};
    $\alpha$ y $CF$ según el Cuadro~\ref{tab:cf-derecho-via}.}
\end{table}

\subsection{Procesamiento con la Cadena de Herramientas}
\label{subsec:anillo-procesamiento}

Las tres herramientas del Capítulo~\ref{ch:Capas-codigo} se usaron con el mismo
código en los tres casos; lo que cambió fue la información que reciben. Cada
escenario corrió en sus propios contenedores Docker, una instancia de
\texttt{Apímetro} y otra de \texttt{VFTModel} con sus datos aislados, de modo
que la red actual y los dos casos con anillo se calcularon por separado, con la
misma versión del modelo y sin que los datos de uno alteraran los de otro.
\texttt{Transport-gis-zmvm-mjg} conserva esa separación al exportar los
resultados de cada caso en carpetas distintas. Esto importa para quien planea:
un urbanista, un arquitecto o un ingeniero puede comparar escenarios sabiendo
que cada uno es independiente y que agregar uno nuevo no compromete la
integridad de los que ya existen. El Cuadro~\ref{tab:anillo-procesamiento}
resume qué hizo cada herramienta en este capítulo y qué entregó.

\begin{table}[H]
    \centering
    \caption[Procesamiento de los casos con anillo]{%
    Papel de cada herramienta en el procesamiento de los Casos~2 y~3.}
    \label{tab:anillo-procesamiento}
    \small
    \begin{tabular}{|p{3.3cm}|p{5.6cm}|p{4.6cm}|}
        \hline
        \textbf{Herramienta} & \textbf{Qué se hizo} & \textbf{Qué entrega} \\
        \hline
        \texttt{Apímetro} (§\ref{sec:apimetro-intro}) &
            Registra las cuatro líneas Periférico Interior y sus estaciones
            como un sistema más de la red &
            Red con anillo para cada caso \\
        \hline
        \texttt{VFTModel} (§\ref{sec:vftmodel-intro}) &
            Construye el grafo de cada caso y calcula los indicadores con la
            corrida 2026-10-04 &
            Indicadores y perfil Garibelt por caso \\
        \hline
        \texttt{Transport-gis-zmvm-mjg} (§\ref{sec:transport-gis}) &
            Cruza los resultados de los tres casos y los lleva a
            \textbf{QGIS} y \textbf{Tableau} &
            Mapas y \termino{dashboards} de cada caso y comparativos \\
        \hline
    \end{tabular}
    \fuentefigura{Elaboración propia.}
\end{table}

\subsection{El Grafo con Anillo}
\label{subsec:anillo-grafo}

Al incorporar el \textsc{api-i}, el grafo crece poco en tamaño. Como se observa
en el Cuadro~\ref{tab:anillo-grafo}, los Casos~2 y~3 suman las mismas 94
estaciones a los 11,115 nodos de la red actual y alrededor de 480 aristas a las
45,517 que ya tenía. Las 94 estaciones quedan dentro de la componente gigante,
que pasa de 10,610 a 10,707 nodos. Los tres nodos restantes de esa diferencia
ya existían: son tres nodos del Metrobús en el norte de la ciudad, entre ellos
Río de los Remedios, que en la red actual estaban desconectados del resto y que
el anillo integra. El anillo, en otras palabras, no cambia la escala del grafo;
lo que puede cambiar es por dónde pasan los caminos más cortos, y eso es lo que
revisan los indicadores de §\ref{sec:anillo-indicadores}.

\begin{table}[H]
    \centering
    \caption[Tamaño del grafo en los tres casos]{%
    Tamaño del grafo multimodal en los tres casos. Corrida: 2026-10-04.}
    \label{tab:anillo-grafo}
    \small
    \begin{tabular}{|l|r|r|r|}
        \hline
        & \textbf{Caso~1} & \textbf{Caso~2} & \textbf{Caso~3} \\
        & \textbf{(red actual)} & \textbf{(Anillo MB)} &
          \textbf{(Anillo Metro)} \\
        \hline
        Nodos                         & 11,115 & 11,209 & 11,209 \\
        Aristas                       & 45,517 & 46,002 & 45,993 \\
        Nodos de la componente gigante & 10,610 & 10,707 & 10,707 \\
        Proporción en la componente gigante & 95.46\,\% & 95.52\,\% &
            95.52\,\% \\
        Aristas de la componente gigante & 44,377 & 44,870 & 44,861 \\
        \hline
    \end{tabular}
    \fuentefigura{Elaboración propia con \texttt{VFTModel}
    \citep{cabrera2026VFTModel}.}
\end{table}
